\documentclass[11pt,a4paper]{article}
\usepackage[T1]{fontenc}
\usepackage{lmodern,amsmath,amssymb,amsthm,microtype}
\usepackage[margin=27mm]{geometry}
\usepackage[colorlinks=true,linkcolor=blue,urlcolor=blue]{hyperref}
\setlength{\emergencystretch}{2em}
\newtheorem{theorem}{Theorem}[section]
\newtheorem{lemma}[theorem]{Lemma}
\newtheorem{corollary}[theorem]{Corollary}
\newcommand{\Q}{\mathbb Q}
\title{Laurent scaling obstructions for two elliptic detectors}
\author{}\date{17 November 2026}
\begin{document}\maketitle
\begin{abstract}
We prove two obstructions to quadratic compilation by diagonal Laurent
scaling with specified unit coordinates. Both use the exact normalization
of the two physical contact weights in the retained elliptic system.
The first concerns a shared special-factor allocation. The second uses
one literal rational reference row and three field factors; it holds
for both the original unit basis and a modified basis of the same size.
The proofs allow arbitrary real weights. Exact integer certificates and
coordinate-map checks require only the Python standard library.
These results concern the stated monomial transformations and give no
general lower bound on the number of variables in quartic definitions.
\end{abstract}

\section{Physical contact weights and retained unit coordinates}

Let $x_i$ denote the coordinates of the table in \cite{Quartic}, and set
\[
                         u=x_{77},\qquad v=x_{90}.
\]
The two physical contact weights are
\begin{equation}\label{eq:physical-weights}
                         q_0=u,\qquad q_1=uv.
\end{equation}
In particular, $v$ is their ratio. The retained inverse coordinates are
$x_{78}=u^{-1}$, $x_{91}=v^{-1}$, and $x_{92}=u^{-1}v^{-2}$.
For formal unit exponents in the basis $(u,v)$, put
\[
\begin{aligned}
\mathcal P_-&=\{(0,0),(1,0),(-1,0),(0,1),(0,-1),(-1,-2)\},\\
\mathcal P_+&=\{(0,0),(1,0),(-1,0),(0,1),(0,-1),(-1,2)\}.
\end{aligned}
\]
The zero vector represents the constant one. The modified basis
$\mathcal P_+$ replaces $x_{92}$ by $y=u^{-1}v^2$, defined by
$uy=v^2$. Since $x_{92}=v^{-4}y$, this is an invertible Laurent
change and adds no coordinate.

The identities \eqref{eq:physical-weights} follow by composing the
retained coordinate maps. The first Laurent reduction sends the
physical $q_1$ coordinate $92$ to coordinate $91$ of the
196-coordinate table. Its parameter-balancing exponent is zero.
The next Laurent map gives this coordinate weight $-1$ in
$w_0=u^{-1}$, and compacts it to coordinate $90$. The final
parameter weight is zero. Thus its physical value is $uv$.
Similarly the physical inverse square is $u^{-1}x_{92}$, giving
$x_{92}=u/q_1^2=u^{-1}v^{-2}$.

We allow a diagonal change $Z=u^{a_Z}v^{b_Z}\widetilde Z$ for every
nonunit coordinate and division of each row by a Laurent monomial.
The weight vector $w_Z=(a_Z,b_Z)$ may even be real. Only literal
monomials are used: no addition of rows or additional polynomial
consequence replaces a nonunit monomial. If a term $u^{e_0}v^{e_1}M$
has nonunit degree $d$, at most $2-d$ retained unit coordinates may
multiply it in a quadratic row. In particular a quadratic nonunit
monomial must have residual unit exponent zero.

\section{A shared special-factor obstruction}

Write $d_0=(1,0)$ and $d_1=(1,1)$ for the physical directions in
\eqref{eq:physical-weights}.

\begin{theorem}\label{thm:special}
Suppose two rows contain the noncancelled terms $q_j^4J^2$ and
$q_j A$, for $j=0,1$, where $J,A$ are shared nonunit coordinates.
Neither $\mathcal P_-$ nor $\mathcal P_+$ permits a quadratic
presentation obtained by the allowed diagonal Laurent transformations.
Common fixed unit factors attached to either displayed term do not
change the conclusion.
\end{theorem}
\begin{proof}
If $r_j$ is the row-division exponent, the square term forces
$r_j=4d_j+2w_J$. The residual exponent of the linear term is
$z_j=w_A-2w_J-3d_j\in\mathcal P_\pm$. Hence
\[
                         z_1-z_0=(0,-3).
\]
In either exponent set, two points with the same first coordinate
have second-coordinate difference at most two. Thus
$(0,-3)\notin\mathcal P_\pm-\mathcal P_\pm$, a contradiction.
Fixed factors common to the two rows cancel on subtraction.
\end{proof}

For example, storing $J=I^4$ in the physical factors
$1+q_j^4 I^8+q_j A+q_j^2B+q_j^3C_1^4$ gives exactly these two
terms. This assertion concerns a shared allocation; distributing
the linear tests differently is outside Theorem~\ref{thm:special}.

\section{A reference-row and factor obstruction}

Consider six nonunit coordinates $v_0,h_1,j_0,j_1,a_0,a_1$ and
the following required support:
\begin{equation}\label{eq:four-support}
\begin{array}{c|l}
\text{row}&\text{noncancelled terms}\\\hline
R&uv_0,\ v_0h_1\\
F_{00}&uh_1,\ u^3j_0^2,\ u^3j_1^2\\
F_{01}&1,\ ua_0^2,\ ua_1^2\\
F_{10}&u^5v^{11}j_0j_1,\ uv a_0a_1.
\end{array}
\end{equation}
Here $R$ is one literal rational reference row. The other three
rows are pairs of field-component equations with common row
multipliers; the justification for this convention is given below.
Coefficients of the displayed monomials need only be nonzero.

\begin{theorem}\label{thm:reference}
Neither $\mathcal P_-$ nor $\mathcal P_+$ permits a quadratic
presentation of \eqref{eq:four-support} by diagonal Laurent
coordinate changes and monomial row division. The conclusion
holds for arbitrary real weights, with no bound on their size.
\end{theorem}
\begin{proof}
Project exponent vectors by $\ell(e_0,e_1)=e_0+e_1$. On either
unit basis the maximum is one. The minimum is $-m$, where $m=3$
for $\mathcal P_-$ and $m=1$ for $\mathcal P_+$.
Write $H,J_i,A_i$ for the projected weights of $h_1,j_i,a_i$, and
$r,r_{00},r_{01},r_{10}$ for the projected row-division weights.

The quadratic term $v_0h_1$ in $R$ determines its row weight.
The residual exponent of the linear term $uv_0$ therefore has
projection $1-H$. Its degree budget is one, so $1-H\le1$ and
$H\ge0$. The linear term $uh_1$ in $F_{00}$ similarly gives
$1+H-r_{00}\le1$. Consequently $r_{00}\ge H\ge0$.

The quadratic terms give the exact equalities
\[
\begin{aligned}
r_{00}&=3+2J_0=3+2J_1=3+J_0+J_1,\\
r_{01}&=1+2A_0=1+2A_1=1+A_0+A_1,\\
r_{10}&=16+J_0+J_1=2+A_0+A_1.
\end{aligned}
\]
It follows that $r_{01}-r_{00}=12$. However, the constant term in
$F_{01}$ has residual exponent $-r_{01}$ and unit degree budget
two. Thus $-r_{01}\ge-2m$, so
\[
                    r_{01}-r_{00}\le2m\le6<12,
\]
a contradiction.
\end{proof}

\subsection{The literal physical interface}

In the table of \cite{Quartic}, take $v_0=x_{19}$ and $h_1=x_{65}$.
The rational row named
\[
\texttt{sharedT.reference.relation3.0}
\]
contains $x_{19}x_{77}$ with coefficient $2128691946$ and
$x_{19}x_{65}$ with coefficient $56549819456$. These are the two
nonzero terms of $R$ used above. The source table is pinned by
SHA-256
\begin{quote}\small\ttfamily\raggedright
304f94843c1a4ac94fcd5f1295fee8b532e7f0424e638314d892c556e2fbfdf9.
\end{quote}

Let $C=149^N$ be the fixed coefficient parameter and $\rho^2=2$.
Composition of all diagonal, affine and parameter maps gives
\[
\begin{aligned}
L_2&=Cu^{-1}(x_{33}+x_{34}\rho),&
L_2^0&=Cu^{-1}(x_{35}+x_{36}\rho),\\
C_2^2&=Cu^{-3}(x_{66}+2x_{67}\rho),&
C_3^2&=Cu^{-3}(x_{68}+2x_{69}\rho),\\
D_H^2&=u^{-1}(x_{64}+2x_{65}\rho),&
B_0^2&=u^{-1}(x_{19}+2x_{20}\rho).
\end{aligned}
\]
For example, the final comparison-square rows are exactly
\[
\begin{aligned}
x_{66}u^{-1}
 &=C\bigl((x_{33}-x_{35})^2+2(x_{34}+x_{36})^2\bigr),\\
x_{67}u^{-1}
 &=C(x_{33}-x_{35})(x_{34}+x_{36}).
\end{aligned}
\]
Thus $j_0=Cx_{66}$ and $j_1=2Cx_{67}$ are diagonal
normalizations with $C_2^2=u^{-3}(j_0+j_1\rho)$.
Similarly write $C_3^2=u^{-3}(k_0+k_1\rho)$ and
$A_1=a_0+a_1\rho$, $B_1=b_0+b_1\rho$. These fixed nonzero
normalizations affect no unit exponent.

The relevant physical factors are
\[
\begin{aligned}
\Phi_{00}&=1+q_0^9C_2^4+q_0^{11}C_3^4+q_0^2D_H^2,\\
\Phi_{01}&=1+q_0A_1^2+q_0^3B_1^2+q_0^2B_0^2,\\
\Phi_{10}&=1+q_1^{11}C_2^4+q_1A_1^2+q_1^2L,
\end{aligned}
\]
where $L$ is a separately stored linear test output. Substituting
$q_0=u$, $q_1=uv$ gives the unit powers $u^5v^{11}$ and $uv$
in $F_{10}$. Expansion into real and imaginary components gives
every factor term in \eqref{eq:four-support}. Subtracting distinct
factor-output coordinates cancels none of these terms.

The retained physical output formulas, including their unit factors,
are
\[
\begin{aligned}
\Phi_{00}&=Cu(x_{118}+x_{119}\rho),\\
\Phi_{01}&=Cu(x_{120}+x_{121}\rho),\\
\Phi_{10}&=Cuv^2(x_{134}+x_{135}\rho).
\end{aligned}
\]
Thus the second detector output has unit exponent $(1,2)$.
Absorbing the fixed nonzero factor $C$ into an output coordinate
does not remove this unit factor. The exact mapping checker verifies
all twenty factor-component output maps, not only these three pairs.

For completeness, common component-row multipliers are forced
even if independent multipliers are initially allowed. If a field
factor contains a full square $u^av^b(Z_0+\rho Z_1)^2$, its real
row contains the nonzero terms $u^av^bZ_0^2$ and
$2u^av^bZ_1^2$. Preserving degree two forces
\[
 (a,b)+2w(Z_0)=r_{\rm real}=(a,b)+2w(Z_1).
\]
Thus the two coordinate weights agree. The nonzero imaginary term
$2u^av^bZ_0Z_1$ then gives $r_{\rm imag}=r_{\rm real}$.
Use $Z=(j_0,j_1)$ in $\Phi_{00},\Phi_{10}$ and
$Z=(a_0,a_1)$ in $\Phi_{01}$. These square monomials belong to
distinct variable blocks and have nonzero coefficients after
collection. The rational reference row has its own independent
multiplier.

Theorem~\ref{thm:reference} uses no other detector row and no
choice of $L$. It excludes the displayed reference/factor
configuration, including arbitrary redistribution of the remaining
linear tests. Other comparison pairings, different reference
equations, row combinations, or identities using other source
relations are outside its hypotheses.

\section{Exact certificates}

The portable specification and checkers are in
\begin{quote}\raggedright
\nolinkurl{negative_results/computation/physical_detector_scaling_v67}.
\end{quote}
The file \texttt{model.json} specifies all retained unit exponents,
the literal reference row, physical field formulas and selected
supports. The file \texttt{farkas\_certificate.json} gives ten
nonnegative integer multipliers for necessary unit-degree
inequalities. Their exact sums are $0\le-12$ for $\mathcal P_-$
and $0\le-20$ for $\mathcal P_+$. These are separate finite
certificates of Theorem~\ref{thm:reference}.

Running \texttt{python check.py --output certificate.json} uses
only the Python standard library. It reconstructs both components
of the physical factors, checks that all selected coefficients are
nonzero, verifies the full-square multiplier witnesses, and checks
the complete integer cancellation. It also checks
$(0,-3)\notin\mathcal P_\pm-\mathcal P_\pm$.
The separate \texttt{check\_mapping.py} composes the source maps,
checks the literal comparison-square equations and verifies the
normalization $q_1=uv$. Both checks pin the full source table when
it is available. No numerical optimizer, SMT solver or bounded
weight search enters either verification.

\begin{thebibliography}{9}\raggedright
\bibitem{Quartic}\emph{A quartic in 195 variables},
\nolinkurl{papers/quartic/quartic195_20261116_v45.tex}; exact table
\nolinkurl{computation/quartic/affine_laurent195_v45/quadratic-family.json}.
\bibitem{Factors}\emph{Five local factors with squared reciprocal tests},
\nolinkurl{papers/quartic/revised_detector_factors_20261117_v66.tex}.
\end{thebibliography}
\end{document}
